% ---- score2 standard preamble (from: python3 score2_tools.py preamble) ----
\documentclass[10pt]{extarticle}
\usepackage[a4paper,margin=1.5cm]{geometry}
\usepackage[T1]{fontenc}
\usepackage{amsmath,amssymb}
\usepackage{graphicx}
\usepackage[export]{adjustbox}
\usepackage{multicol}
\usepackage{enumitem}
\usepackage{booktabs,array,tabularx}
\setlength{\columnsep}{0.7cm}
\setlength{\parindent}{0pt}
\setlength{\parskip}{3pt}
\pagestyle{empty}
\makeatletter
% figure registry (figdefs.tex is generated by the extract step)
\newcommand{\qfigdef}[4]{\@namedef{qf@file@#1}{#2}\@namedef{qf@w@#1}{#3}\@namedef{qf@wide@#1}{#4}}
\newcommand{\qf@check}[1]{\@ifundefined{qf@file@#1}{\PackageError{score2}{Unknown figure id #1}{}}{}}
\newcommand{\qf@img}[4][]{\includegraphics[width=#2\dimexpr\@nameuse{qf@w@#4}pt\relax,max width=#3,max height=0.45\textheight,#1]{\@nameuse{qf@file@#4}}}
% Full-width questions: finalize wraps every question that contains a wide figure in \qfullbegin...\qfullend,
% so the WHOLE question (stem, figures, options) is set across the page in reading order, then the two columns
% resume. The question's head - its number, stem and first figure (\qfullsplit marks the end of the head) -
% is kept on one page, so a number is never left at a page foot with its figure overleaf; the rest may break
% across pages normally, so no big gaps are left. Without a \qfullsplit, the whole question is kept together
% if it is under 45% of a page. Never write these yourself.
\newif\ifqf@full
\newif\ifqf@boxed
\newsavebox{\qf@box}
\newcommand{\qf@open}{\global\qf@fulltrue\par\global\qf@boxedtrue
  \setbox\qf@box\vbox\bgroup\hsize\textwidth\linewidth\textwidth}
\newcommand{\qf@emit}[1]{\par\egroup\global\qf@boxedfalse
  \ifdim\dimexpr\ht\qf@box+\dp\qf@box\relax<#1\textheight\noindent\usebox{\qf@box}\par\else\unvbox\qf@box\fi}
\newcommand{\qfullsplit}{\ifqf@boxed\qf@emit{0.85}\fi}
\newcommand{\qf@close}{\ifqf@boxed\qf@emit{0.45}\fi\global\qf@fullfalse}
\newcommand{\qfullbegin}{\end{multicols}\qf@open}
\newcommand{\qfullend}{\qf@close\begin{multicols}{2}}
\newcommand{\qfullstart}{\qf@open}
\newcommand{\qfulldone}{\qf@close}
% \qfig[scale]{id}: block figure, centred, directly where it is written (size set by the tool).
% Fallback if a wide figure is not inside a full-width question: break out of the columns in place - no floats.
\newcommand{\qfig}[2][1]{\qf@check{#2}%
  \ifqf@full\qf@blk{#1}{#2}\else
  \ifnum\@nameuse{qf@wide@#2}=1 \expandafter\@firstoftwo\else\expandafter\@secondoftwo\fi
  {\end{multicols}\begin{center}\qf@img{#1}{\textwidth}{#2}\end{center}\begin{multicols}{2}}%
  {\qf@blk{#1}{#2}}\fi}
\newcommand{\qf@blk}[2]{\par\smallskip{\centering\qf@img{#1}{\linewidth}{#2}\par}\smallskip}
% \qfigi[scale]{id}: inline figure (inside option lists / tables / minipages); never breaks columns.
\newcommand{\qfigi}[2][1]{\qf@check{#2}\qf@img[valign=c]{#1}{\linewidth}{#2}}
\makeatother
% structure macros - lint reads these to check completeness and figure placement
\newcommand{\papersec}[1]{\par\bigskip{\centering\Large\bfseries PAPER-#1\par}\medskip}
\newcommand{\partsec}[2]{\par\medskip\noindent{\large\bfseries PART-#1: #2}\par\smallskip}
\newcommand{\sect}[1]{\par\smallskip\noindent\textbf{#1}\par}
\newcommand{\qn}[1]{\par\smallskip\noindent\textbf{#1.}~}
\newcommand{\paraq}[2]{\par\smallskip\noindent\textbf{Paragraph for Questions #1 and #2}\par}
\newenvironment{qoptions}{\begin{itemize}[leftmargin=1.8em,itemsep=1pt,topsep=2pt,label={}]}{\end{itemize}}
\newcommand{\opt}[2]{\item[(#1)] #2}
% figure registry (generated by score2_tools.py finalize)
\qfigdef{p1_phy_s1_q01_1}{Pre-Score (2023-24) - Test 01 (08 Oct 2023)-img/p1_phy_s1_q01_1.png}{134.9}{0}
\qfigdef{p1_phy_s1_q03_1}{Pre-Score (2023-24) - Test 01 (08 Oct 2023)-img/p1_phy_s1_q03_1.png}{153.8}{0}
\qfigdef{p1_phy_s1_q04_1}{Pre-Score (2023-24) - Test 01 (08 Oct 2023)-img/p1_phy_s1_q04_1.png}{108.2}{0}
\qfigdef{p1_phy_s2_q05_1}{Pre-Score (2023-24) - Test 01 (08 Oct 2023)-img/p1_phy_s2_q05_1.png}{123.2}{0}
\qfigdef{p1_phy_s2_q05_2}{Pre-Score (2023-24) - Test 01 (08 Oct 2023)-img/p1_phy_s2_q05_2.png}{125.9}{0}
\qfigdef{p1_phy_s2_q06_1}{Pre-Score (2023-24) - Test 01 (08 Oct 2023)-img/p1_phy_s2_q06_1.png}{103.5}{0}
\qfigdef{p1_phy_s3_para1_2_1}{Pre-Score (2023-24) - Test 01 (08 Oct 2023)-img/p1_phy_s3_para1_2_1.png}{224.4}{0}
\qfigdef{p1_phy_s3_para3_4_1}{Pre-Score (2023-24) - Test 01 (08 Oct 2023)-img/p1_phy_s3_para3_4_1.png}{260.4}{1}
\qfigdef{p1_phy_s4_q07_1}{Pre-Score (2023-24) - Test 01 (08 Oct 2023)-img/p1_phy_s4_q07_1.png}{245.2}{0}
\qfigdef{p1_phy_s4_q09_1}{Pre-Score (2023-24) - Test 01 (08 Oct 2023)-img/p1_phy_s4_q09_1.png}{106.6}{0}
\qfigdef{p1_math_s1_q03_1}{Pre-Score (2023-24) - Test 01 (08 Oct 2023)-img/p1_math_s1_q03_1.png}{88.4}{0}
\qfigdef{p1_math_s3_para5_6_1}{Pre-Score (2023-24) - Test 01 (08 Oct 2023)-img/p1_math_s3_para5_6_1.png}{96.5}{0}
\qfigdef{p2_phy_s1_q02_1}{Pre-Score (2023-24) - Test 01 (08 Oct 2023)-img/p2_phy_s1_q02_1.png}{119.0}{0}
\qfigdef{p2_phy_s1_q03_1}{Pre-Score (2023-24) - Test 01 (08 Oct 2023)-img/p2_phy_s1_q03_1.png}{268.6}{1}
\qfigdef{p2_phy_s1_q04_1}{Pre-Score (2023-24) - Test 01 (08 Oct 2023)-img/p2_phy_s1_q04_1.png}{227.0}{0}
\qfigdef{p2_phy_s1_q05_1}{Pre-Score (2023-24) - Test 01 (08 Oct 2023)-img/p2_phy_s1_q05_1.png}{92.1}{0}
\qfigdef{p2_phy_s1_q06_1}{Pre-Score (2023-24) - Test 01 (08 Oct 2023)-img/p2_phy_s1_q06_1.png}{153.0}{0}
\qfigdef{p2_phy_s2_para7_8_1}{Pre-Score (2023-24) - Test 01 (08 Oct 2023)-img/p2_phy_s2_para7_8_1.png}{66.4}{0}
\qfigdef{p2_phy_s2_q07_A}{Pre-Score (2023-24) - Test 01 (08 Oct 2023)-img/p2_phy_s2_q07_A.png}{78.5}{0}
\qfigdef{p2_phy_s2_q07_B}{Pre-Score (2023-24) - Test 01 (08 Oct 2023)-img/p2_phy_s2_q07_B.png}{74.0}{0}
\qfigdef{p2_phy_s2_q07_C}{Pre-Score (2023-24) - Test 01 (08 Oct 2023)-img/p2_phy_s2_q07_C.png}{75.4}{0}
\qfigdef{p2_phy_s2_q07_D}{Pre-Score (2023-24) - Test 01 (08 Oct 2023)-img/p2_phy_s2_q07_D.png}{63.0}{0}
\qfigdef{p2_phy_s2_para9_10_1}{Pre-Score (2023-24) - Test 01 (08 Oct 2023)-img/p2_phy_s2_para9_10_1.png}{108.5}{0}
\qfigdef{p2_phy_s3_para1_2_1}{Pre-Score (2023-24) - Test 01 (08 Oct 2023)-img/p2_phy_s3_para1_2_1.png}{208.6}{0}
\qfigdef{p2_phy_s3_para3_4_1}{Pre-Score (2023-24) - Test 01 (08 Oct 2023)-img/p2_phy_s3_para3_4_1.png}{108.5}{0}
\qfigdef{p2_phy_s3_para5_6_1}{Pre-Score (2023-24) - Test 01 (08 Oct 2023)-img/p2_phy_s3_para5_6_1.png}{92.6}{0}
\qfigdef{p2_phy_s4_q07_1}{Pre-Score (2023-24) - Test 01 (08 Oct 2023)-img/p2_phy_s4_q07_1.png}{99.7}{0}
\qfigdef{p2_chem_s1_q01_1}{Pre-Score (2023-24) - Test 01 (08 Oct 2023)-img/p2_chem_s1_q01_1.png}{187.5}{0}
% ---- body skeleton ----
% \begin{document}
% \begin{center}{\Large\bfseries JEE (Main+Advanced) Score-II --- Test NN (DD Mon YYYY)}\end{center}
% \begin{multicols}{2}
%   \papersec{1}   (only for two-paper tests)
%   \partsec{1}{PHYSICS} \sect{Section-I (i): Single Correct} \qn{1} ... \qfig{p1_phy_q01_1} ...
% \end{multicols}
% \end{document}

\begin{document}
\begin{center}{\Large\bfseries Pre-Score (2023-24) --- Test 01 (08 Oct 2023)}\end{center}
\begin{multicols}{2}
\papersec{1}
\partsec{1}{PHYSICS}
\sect{SECTION-I(i)}
\qn{1} A bug is moving with velocity $(2\hat i-3\hat j+4\hat k)$ m/s in front of plane mirror which is also moving with velocity $(2\hat i+\hat j+\hat k)$ m/s. Plane mirror always remain parallel to y-z plane. Velocity of image of bug is :-
\qfig{p1_phy_s1_q01_1}
\begin{qoptions}
\opt{A}{$2\hat i-3\hat j+4\hat k$}
\opt{B}{$-2\hat i-3\hat j+4\hat k$}
\opt{C}{$-2\hat i+3\hat j-4\hat k$}
\opt{D}{None}
\end{qoptions}
\qn{2} Ram is looking at his face in a mirror kept 10 cm away and he finds that his image is erect and magnified ($m=1.8$). If he holds the mirror 50 cm away :
\begin{qoptions}
\opt{A}{He cannot see the image because reflected rays falling on his eyes are converging}
\opt{B}{He sees a magnified and erect image}
\opt{C}{He sees a diminished and inverted image}
\opt{D}{He sees a magnified and inverted image.}
\end{qoptions}
\qn{3} One end of a glass rod of refractive index $n=1.5$ is a spherical surface of radius of curvature $R$. The centre of the spherical surface lies inside the glass. A point object placed in air on the axis of the rod at the point $P$ has its real image inside glass at the point $Q$ (see fig.). A line joining the points $P$ and $Q$ cuts the surface at $O$ such that $OP=2OQ$. The distance $PO$ is :-
\qfig{p1_phy_s1_q03_1}
\begin{qoptions}
\opt{A}{$8R$}
\opt{B}{$7R$}
\opt{C}{$2R$}
\opt{D}{None of these}
\end{qoptions}
\qn{4} Radii of curvature of surfaces of a biconvex lens are 10 cm and 20 cm and refractive index of its material is $\mu=1.5$. If the refractive indices of the medium on the two sides of the lens are $\mu_1$ and $\mu_2$, then the parallel rays incident on the lens from left will focus at a distance of $f_2$ (distance being measured from the optical centre of the lens) where $f_2$ is equal to :-
\qfig{p1_phy_s1_q04_1}
\begin{qoptions}
\opt{A}{$\dfrac{20\mu_1}{4.5-(2\mu_1+\mu_2)}$}
\opt{B}{$\dfrac{20\mu_2}{4.5-(\mu_1+2\mu_2)}$}
\opt{C}{$\dfrac{20\mu_2}{4.5-(2\mu_1+\mu_2)}$}
\opt{D}{$\dfrac{20\mu_1}{4.5-(\mu_1+2\mu_2)}$}
\end{qoptions}
\sect{SECTION-I(ii)}
\qn{5} A monochromatic light is incident on a prism as shown. Deviation suffered by ray varies with angle of incidence as shown. Assuming refraction takes place at 1 \& 2. Choose the correct alternative(s).
\qfig{p1_phy_s2_q05_1}
\qfig{p1_phy_s2_q05_2}
\begin{qoptions}
\opt{A}{$i_1$ can be equal to zero if $A=30^\circ$}
\opt{B}{If angle of incidence is increased by $2^\circ$, then angle of emergence may decrease by less than $2^\circ$.}
\opt{C}{If angle of incidence is increased by $2^\circ$, then angle of emergence may decrease by more than $2^\circ$.}
\opt{D}{If angle of incidence is increased by $2^\circ$, then angle of emergence may increase by $2^\circ$.}
\end{qoptions}
\qn{6} A small source of light is mounted inside a cylindrical container of height $h$. The bottom of the container is covered with a mirror. Initially, the container is empty. Then a clear liquid with the index of refraction $n$ is slowly poured into the container. The level of liquid $H$ rises steadily, reaching the top of the container in time $T$. Let $h_1$ be the distance of source of the light from the bottom of the container. Consider paraxial ray approximation. The observer is vertically above the source in air and he is observing image after reflection.
\qfig{p1_phy_s2_q06_1}
Consider two cases\\
(1) $H\le h_1$\quad (2) $h_1\le H\le h$
\begin{qoptions}
\opt{A}{The speed of the image of the source in case (1) during this process is $\dfrac{2h}{T}\Bigl(1-\dfrac1n\Bigr)$.}
\opt{B}{The speed of the image of the source in case (2) during this process is $\dfrac hT\Bigl(1-\dfrac1n\Bigr)$.}
\opt{C}{The speed of the image of the source in case (1) during this process is $\dfrac h{2T}\Bigl(1-\dfrac1n\Bigr)$.}
\opt{D}{The speed of the image of the source in case (2) during this process is $\dfrac h{2T}\Bigl(\dfrac1n\Bigr)$.}
\end{qoptions}
\qn{7} Consider an equilateral prism of refractive index $\mu$. It is found that when light is incident at an angle $60^\circ$, angle of emergence is `$e$', while when light is incident at an angle $50^\circ$, it undergoes TIR on the 2nd refracting surface. Choose the CORRECT options :-
\begin{qoptions}
\opt{A}{A ray incident at an angle $65^\circ$ will surely emerge from other surface.}
\opt{B}{A ray incident at an angle of $40^\circ$ will surely suffer TIR on other surface.}
\opt{C}{A ray incident at an angle of $50.7^\circ$, will surely suffers TIR on other surface.}
\opt{D}{A ray incident at an angle $59^\circ$, may suffer TIR on the other surface.}
\end{qoptions}
\qn{8} Two mirrors are inclined at an angle of $30^\circ$ one of the ray is incident parallel to the one of the mirror. Then select the correct statement(s)-
\begin{qoptions}
\opt{A}{Ray will undergo five reflections from the system}
\opt{B}{Total deviation from the system will $120^\circ$}
\opt{C}{Ray will retrace it path finally}
\opt{D}{Total deviation suffered by the ray will be $90^\circ$}
\end{qoptions}
\qn{9} A point source of light is located on the optical axis of a thin converging lens of aperture diameter 1 cm with focal length 10 cm. After the refraction from the lens
\begin{qoptions}
\opt{A}{The maximum angle of diverging beam possible is 0.1 radians.}
\opt{B}{The maximum angle of a converging beam possible is 0.1 radians.}
\opt{C}{The minimum possible distance between the real image and the object is 40 cm.}
\opt{D}{The maximum possible angle of deviation for any ray is 0.05 radians.}
\end{qoptions}
\qn{10} An observer holds in front of himself a thin equi convex converging lens; $a$ is the numerical value of the radius of curvature of each face. He sees two clear and distinct images of one of his eyes, one erect and the other inverted. The refractive index of lens is 1.50. Least distance of distinct vision is 25 cm. One image is due to reflection from front surface \& other is due to reflection from rear surface.
\begin{qoptions}
\opt{A}{The observer must be at a distance more than $0.5a$ from the lens.}
\opt{B}{The distance between observer and his inverted image must be more than 25 cm.}
\opt{C}{If the observer comes at a distance more than $0.25a$ but less than $0.5a$ from the lens, he sees two erect images.}
\opt{D}{If the observer comes at a distance less than $0.25a$ from the lens, he can see two erect images.}
\end{qoptions}
\sect{SECTION-II(i)}
\paraq{1}{2}
For the given system of blocks and a concave mirror answer the following. (Assuming $A$, $B$ \& $C$ are very small). Focal length of mirror is 20 cm height of $A$ \& $B$ are same.
\qfig{p1_phy_s3_para1_2_1}
\qn{1} If $A$ is at $x=10$ cm then image of $B$ and image of $A$ are of same size. What is the distance of $B$ (in cm) from mirror ?
\qn{2} Velocity of image of $A$ when it is at $x=30$ cm is\underline{\hspace{1.5cm}}m/s
\qfullbegin
\paraq{3}{4}
An experiment is performed to calculate the focal length of a convex mirror using convex lens. The setup consists of two pins, optical bench, one convex lens and one convex mirror. When only convex lens is placed at a distance 60 cm from pin (object), its image is formed at a distance 30 cm from the lens as shown in figure-(i). Now convex mirror is also kept right to the lens without changing distance between pin (object) and lens. The mirror is kept at a distance 10 cm from lens such that position of image of object pin and position of object pin coincides as shown in figure-(ii)
\qfig{p1_phy_s3_para3_4_1}
\qfullend
\qn{3} Focal length of convex lens used (in cm) is :
\qn{4} Focal length of convex mirror (in cm) is :
\paraq{5}{6}
Assume that each grade of glass is characterized by a refractive index $n_0$ at a wavelength $\lambda_0=0.55\ \mu$m and a coefficient dispersion $\beta$, which shows how much the refractive index of a given type of glass depends on wavelength $\lambda$. $n_\lambda=n_0+\beta(\lambda-\lambda_0)$. There are two grades of clear glasses with characteristics $\beta_1\ne\beta_2$. We want to make lenses for glasses that lack chromatic aberration, i.e. combined optical power is not dependent on the wavelength. For this we need to combine two or more lens. Assume that we have combined a bi-convex lens with a convexo-concave lens. Radii of bi-convex lens is $R_1$ and $R_2$ and that for convexo-concave lens is $R_2$ and $R_0$. The lens exactly fit into each other. $\beta_1=0.05\ \mathrm{nm}^{-1}$, $\beta_2=0.06\ \mathrm{nm}^{-1}$, $R_0=20$ cm, $R_1=10$ cm, $n_{01}=n_{02}=1.5$
\qn{5} What is the value of $R_2$ (in cm)?
\qn{6} If $\beta_1=0.05\ \mathrm{nm}^{-1}$, $\beta_2=0.06\ \mathrm{nm}^{-1}$, $R_0=20$ cm, $R_1=10$ cm, $n_{01}=1.5$ and $n_{02}=1.6$ $\lambda_0=500$ nm. What is the value of $R_2$ (in cm) for which the ray of wavelength 520 nm passes through the combination without deviation?
\sect{SECTION-II(ii)}
\qn{7} A curved thick glass surface is silvered at curved face \& not silvered on plane surface. Object is placed at $A$ as shown in figure. Considering $P$ (pole of the silvered surface) as origin. If $x$ - co-ordinate of final image is `$n$' cm then find $n$.
\qfig{p1_phy_s4_q07_1}
\qn{8} A flint glass convex lens of focal length 16 cm is placed in contact with crown glass lens. The refractive indices of crown glass for violet and red colour are $\mu_v=1.525$, $\mu_r=1.515$ and for flint glass $\mu_v'=1.655$ and $\mu_r'=1.645$. This forms an achromatic combination with the flint glass lens. What is the focal length of the combination (in cm) ?
\qn{9} The concave spherical mirror of radius of $R=9$ cm forms the image of a source which coincides with a source. When a little liquid is poured in the mirror, a second image is formed between the source and the mirror at a distance $\ell=4$ cm from the source. Find the refractive index of the liquid. Write `$5\mu$' in OMR sheet.
\qfig{p1_phy_s4_q09_1}
\partsec{2}{CHEMISTRY}
\sect{SECTION-I(i)}
\qn{1} Statement-1 : $\Delta_iH$ of H is 1312 kJ/mol.\\
Statement-2 : Reactivity of hydrogen is very low as compared to halogens.\\
Statement-3 : Deuterium exist mostly in terms of HD.\\
Statement-4 : Tritium is radioactive and emits high energy $\beta$ particles.\\
Statement-5 : H--H bond dissociation enthalpy is lowest for a single bond.\\
Consider above statements and select correct scheme (T = True, F = False)
\begin{qoptions}
\opt{A}{T T F T F}
\opt{B}{T T T T F}
\opt{C}{F T T T T}
\opt{D}{T T T F F}
\end{qoptions}
\qn{2} Select correct statement :
\begin{qoptions}
\opt{A}{$\alpha$-nitroso-$\beta$-napthol is chelating ligand.}
\opt{B}{Cupron is ligand with only N-donor sites.}
\opt{C}{D$_2$O is used in production of cephalosporin.}
\opt{D}{Certain plants like Pinus, Juniparus, Pyrus can metabolise sulphur oxides.}
\end{qoptions}
\qn{3} Select the incorrect statement about ice :
\begin{qoptions}
\opt{A}{At atmospheric pressure ice crystallises in the Hexagonal form.}
\opt{B}{At very low temperatures ice condenses to Cubic form.}
\opt{C}{Ice formed by heavy water sinks in normal water.}
\opt{D}{In ice each oxygen atom is surrounded tetrahedrally by four other oxygen at two type of distance 95.7 pm and 276 pm.}
\end{qoptions}
\end{multicols}
\noindent\begin{minipage}{\textwidth}\qn{4}\par\vspace{-3pt}
\begin{center}\small
\begin{tabular}{|c|l|c|p{6.2cm}|}
\hline
1. & Ni(CO)$_4$ & (a) & A complex ion having a magnetic moment of 2.83 Bohr magnetons \\\hline
2. & [Ni(CN)$_4$]$^{2-}$ & (b) & 34 \\\hline
3. & [Ni(NH$_3$)$_6$]$^{2+}$ & (c) & Ziese's salt \\\hline
4. & [Zn(NH$_3$)$_4$]$^{2+}$ & (d) & cis-trans Isomerism \\\hline
5. & [Pd(NH$_3$)$_2$Cl$_2$] & (e) & 54 \\\hline
6. & K$^+$[Pt(C$_2$H$_4$)Cl$_3$]$^-$ & (f) & A low-spin (spin-paired) complex \\\hline
7. & [Co(NH$_3$)$_4$(H$_2$O)$_2$]Cl & (g) & A complex ion having a magnetic moment of 1.73 Bohr magnetons \\\hline
8. & EAN for [Ni(CN)$_4$]$^{2-}$ & (h) & Linear complex \\\hline
9. & EAN for [Pd(NH$_3$)$_6$]$^{4+}$ & (i) & Shows geometrical isomerism \\\hline
10. & [Co(NH$_3$)$_6$]$^{3+}$ & (j) & A diamagnetic species \\\hline
11. & CoF$_6^{3-}$ & (k) & Square planar \\\hline
12. & [Mn(CN)$_6$]$^{4-}$ & (l) & A high-spin (spin-free) complex \\\hline
13. & [V(H$_2$O)$_6$]$^{3+}$ & (m) & Tetrahedral \\\hline
14. & [Ag(NH$_3$)$_2$]Cl & (n) & A paramagnetic species \\\hline
\end{tabular}
\end{center}\end{minipage}
\begin{multicols}{2}
Consider above given list.\\
Select option in which only correct matches given :
\begin{qoptions}
\opt{A}{1-m, 2-k, 3-h, 8-b, 9-e, 10-d}
\opt{B}{3-h, 4-j, 5-i, 11-l, 12-g, 13-j}
\opt{C}{13-a, 14-h, 9-e, 10-f, 3-j, 4-j}
\opt{D}{7-d, 8-b, 2-k, 3-n, 12-g, 6-c}
\end{qoptions}
\sect{SECTION-I(ii)}
\qn{5} Select the correct statement(s) about actinoids :
\begin{qoptions}
\opt{A}{Actinoids gives mixture of oxides and hydrides on boiling with water.}
\opt{B}{Most of actinoids are slightly affected by nitric acid but react with HCl.}
\opt{C}{Shows voilent reaction with NaOH.}
\opt{D}{Lanthanoid contraction is greater from element to element than actinoid contraction.}
\end{qoptions}
\qn{6} Symmetrical Di-$\mu$-\allowbreak hydroxytetrakis\allowbreak (ethylenediamine)\allowbreak dicobalt(III) chloride (P) complex reacts with aq. HCl to give a product Co(en)$_2$Cl$_3$(Q) which is resolvable with optical isomers. When it is kept in acidic medium the compound changes color and give isomer (R) which is not optically resolvable.\\
Select correct statement(s) ?
\begin{qoptions}
\opt{A}{If one `en' is replaced by two `NH$_3$' molecules then complex (R) become optically resolvable.}
\opt{B}{If one `en' is replaced by two `NH$_3$' molecules and second `en' is replaced by two H$_2$O molecules. Complex `R' still optically not resolvable.}
\opt{C}{If one `en' is replaced by two `NH$_3$' molecules in complex `Q', complex is still optically resolvable.}
\opt{D}{If one `en' is replaced by two `NH$_3$' then second `en' is replaced by two `H$_2$O' in complex `Q' become optically not resolvable.}
\end{qoptions}
\qn{7} A compound (X) is made by mixing Cobalt(III) nitrite and potassium nitrite solution in ratio of 1 : 3. The aqueous solution of the compound showed 4 particles per molecule where as molar conductivity reveals presence of six electrical charges.\\
Select correct statement about complex (X) :
\begin{qoptions}
\opt{A}{Complex `X' is low spin complex.}
\opt{B}{Complex `X' can show maximum `10' isomeric forms.}
\opt{C}{Complex `X' can show `3' optically active isomers.}
\opt{D}{Complex `X' is inner orbital complex.}
\end{qoptions}
\qn{8} The Co(III) aqua complex slowly undergoes redox decomposition to give the Co(II) aqua complex and oxygen, but in presence of `en', Co(II) is slowly oxidised to [Co(en)$_3$]$^{3+}$ by aerial oxygen.\\
Select correct statement(s) :
\begin{qoptions}
\opt{A}{After redox decomposition magnetic moment decreases.}
\opt{B}{After arial oxidation magnetic moment increased.}
\opt{C}{After arial oxidation crystal field splitting energy increased.}
\opt{D}{After redox decomposition crystal field splitting energy decreased.}
\end{qoptions}
\qn{9} Select the correct statement(s) ?
\begin{qoptions}
\opt{A}{[Co(gly)$_2$(OH$_2$)$_2$]$^+$ can produce 5 geometrical isomers.}
\opt{B}{[Tris(tetramine-$\mu$-\allowbreak dihydroxido\allowbreak cobalt(III))]\allowbreak cobalt(III) ion shows optical activity.}
\opt{C}{[PtCl$_2$(NH$_3$)$_2$(NO$_2$)$_2$] can have maximum 6 isomers.}
\opt{D}{For [Co(Br)(en)(NH$_3$)$_2$(OH$_2$)]$^{2+}$ ion total geometrical isomers is equal to the number of optically active isomers.}
\end{qoptions}
\qn{10} Select the incorrect statement(s) :
\begin{qoptions}
\opt{A}{[Cd(CH$_3$NH$_2$)$_4$]$^{2+}$ is less stable as compared to [Cd(en)$_2$]$^{2+}$ due to entropy of system is more in case of [Cd(en)$_2$]$^{2+}$.}
\opt{B}{[Ni(en)$_3$]$^{2+}$ is less stable than [Ni(NH$_3$)$_6$]$^{2+}$}
\opt{C}{[Co(CN)$_6$]$^{3-}$ is more stable than [Rh(CN)$_6$]$^{3-}$}
\opt{D}{[Ag(NH$_3$)$_2$]$^+$ is more stable than [Ag(en)]$^+$}
\end{qoptions}
\sect{SECTION-II(i)}
\paraq{1}{2}
As we know that water is an excellent solvent, water soluble inorganic chemicals that include heavy metals such as cadmium, mercury, nickel etc. constitute an important class of pollutants. All these metals are dangerous to humans because our body cannot excrete them. Over the time, it crosses the tolerance limit. These metals then can damage kidneys, central nervous system, liver etc. Acids (like sulphuric acid) from mine drainage and salts from many different sources including raw salt used to melt snow and ice in the colder climates (sodium and calcium chloride) are water soluble chemical pollutants.
\qn{1} The maximum prescribed limit of nitrate ion in drinking water is `X' mg dm$^{-3}$ and lead is $Y\times10^{-3}$ mg dm$^{-3}$, then calculate $(X+Y)$.
\qn{2} The maximum prescribed limit of Fe metal in drinking water is $P$ mg dm$^{-3}$ and Mn is $Q$ mg dm$^{-3}$. Calculate $(P+Q)$.
\paraq{3}{4}
Two complexes (A) and (B) of Cobalt have same molecular formula Co(NH$_3$)$_3$(H$_2$O)$_2$ClBr$_2$. 34.15 gm of (A) loses 1.8 gm of weight when heated to 200\textdegree C whereas no weight loss was observed on similar treatment with (B).\\
An aqueous solution of (A) has conductivity equivalent to that of a compound with two ions per formula unit while conductivity of (B) in aqueous solution is equivalent to a compound with (X) number ions per unit. Also a 100 ml 0.1 M aq. solution of (A), on treatment with excess of AgNO$_3$. Solution produce (Y) gm of only pale yellow ppt while a 100 ml 0.1 M aqueous solution of (B) on treatment with excess of AgNO$_3$ solution product 3.315 gm precipitate.\\
{[}Molar mass : H = 1, N = 14, O = 16, Cl = 35.5, Co = 59, Br = 80, Ag = 108{]}
\qn{3} the value of (X) is :
\qn{4} The value of Y is :
\paraq{5}{6}
Different theories like electrostatic theory (developed by Kossel, Magnus, Van Arekel, Garrick), Sidgwick's effective atomic number rule (EAN), Pauling's valence bond theory (VBT), crystal field theory (CFT), adjusted crystal field theory (ACFT) or ligand field theory (LFT) providing some allowance of covalence in CFT, and molecular orbital theory (MOT) have been proposed for understanding of bonding in coordination compounds.
\qn{5} How many complexes are low spin type ?\\
$[\mathrm{Co(NH_3)_6}]^{+3}$, $[\mathrm{CoF_6}]^{-3}$, $[\mathrm{Fe(CN)_6}]^{-4}$, $[\mathrm{Fe(H_2O)_6}]^{+2}$, $[\mathrm{Cr(gly)_3}]^{0}$, $[\mathrm{CoBr_2Cl_2}]^{-2}$
\qn{6} Calculate value of $(x+y)$ :\\
Here\\
Among the species given below, the total number of diamagnetic species is $x$.\\
MnCl$_4^{2-}$, NO$_2$ monomer, O$_2^-$ (superoxide), CoCl$_4^{2-}$, Mn$_3$O$_4$, (NH$_4$)$_2$[FeCl$_4$], CuCl$_4^{2-}$, K$_2$MnO$_4$, K$_2$CrO$_4$\\
The total number of complexes which are following Sidgwick EAN rule is $y$.\\
(Atomic number : Mo $\Rightarrow$ 42, Pt $\Rightarrow$ 78, Ag = 47)\\
$[\mathrm{Mo(\sigma\text{--}C_3H_5)Br(NH_3)_2}]$, $[\mathrm{Fe(\eta^5\text{--}C_5H_5)_2}]$, $[\mathrm{Ni(DMG)_2}]$, $[\mathrm{Mn_2(CO)_{10}}]$, $[\mathrm{Fe(CO)_2(NO)_2}]^{0}$, $[\mathrm{PtCl_3(\eta^2\text{--}C_2H_4)}]^{\ominus}$, $[\mathrm{Ag(CN)_2}]^{-1}$, $[\mathrm{Fe_2(CO)_9}]^{0}$
\sect{SECTION-II(ii)}
\qn{7} Calculate the value of $(P\times Q\times R)$\\
Here\\
$P$ = Number of green house gases in the following given set of species.\\
(CO$_2$, O$_2$, N$_2$O, CH$_4$, CFC, CO, H$_2$O(g), O$_3$)\\
$Q$ = Number of nonviable particulate pollutants in the following given set of species.\\
(Smoke particles, Fungi, Fly ash, Moulds, Saw dust, Bacteria, Pulverized coke.)\\
$R$ = Number of Tropospheric gaseous pollutants in the following given set of species.\\
(SO$_2$, Mist, CO$_2$, O$_2$, NO$_2$, Smog, CH$_4$, CO)
\qn{8} Consider the following reactions :\\
(a) H$_2$O$_2$ + HIO$_4$ $\longrightarrow$\\
(b) H$_2$O$_2$ + MnO$_4^-$ $\xrightarrow{\mathrm{H^+}}$\\
(c) H$_2$O$_2$ + Fe$^{2+}$ $\xrightarrow{\mathrm{H^+}}$\\
(d) H$_2$O$_2$ + PbS $\xrightarrow{\mathrm{H^+}}$\\
(e) H$_2$O$_2$ + Ce$^{4+}$ $\longrightarrow$\\
(f) H$_2$O$_2$ + SO$_3^{2-}$ $\longrightarrow$\\
(g) H$_2$O$_2$ + Ag$_2$O $\longrightarrow$\\
(h) H$_2$O$_2$ + [Fe(CN)$_6$]$^{3-}$ $\longrightarrow$\\
(i) H$_2$O$_2$ + Au$^{+3}$ $\longrightarrow$\\
(j) H$_2$O$_2$ + I$_2$ $\xrightarrow{\mathrm{OH^-}}$\\
(k) H$_2$O$_2$ + MnCl$_2$ $\xrightarrow{\mathrm{OH^-}}$\\
(l) H$_2$O$_2$ + HOCl $\xrightarrow{\mathrm{H^+}}$\\
If total number of reactions in which H$_2$O$_2$ act as reducing agent $=X$.\\
If total number of reactions in which H$_2$O$_2$ act as oxidizing agent $=Y$.\\
Then find the value of $X\cdot Y$
\qn{9} Find the number pair of ions which have same outermost E.C.\\
\textbf{Set I :} (Ac$^{3+}$, Th$^{4+}$), (Th$^{3+}$, Pa$^{4+}$), (Pa$^{3+}$, U$^{4+}$), (U$^{3+}$, Np$^{4+}$), (Np$^{3+}$, Pu$^{4+}$), (Pu$^{3+}$, Am$^{4+}$), (Am$^{3+}$, Cm$^{4+}$), (Cm$^{3+}$, Bk$^{4+}$), (Bk$^{3+}$, Cf$^{4+}$), (Cf$^{3+}$, Es$^{4+}$), (Es$^{3+}$, Fm$^{4+}$), (Fm$^{3+}$, Md$^{4+}$), (Md$^{3+}$, No$^{4+}$), (No$^{3+}$, Lr$^{4+}$)\\
\textbf{Set II :} (La$^{2+}$, Ce$^{3+}$), (Ce$^{2+}$, Pr$^{3+}$), (Pr$^{2+}$, Nd$^{3+}$), (Nd$^{2+}$, Pm$^{3+}$), (Pm$^{2+}$, Sm$^{3+}$), (Sm$^{2+}$, Eu$^{3+}$), (Eu$^{2+}$, Gd$^{3+}$), (Gd$^{2+}$, Tb$^{3+}$), (Tb$^{2+}$, Dy$^{3+}$), (Dy$^{2+}$, Ho$^{3+}$), (Ho$^{2+}$, Er$^{3+}$), (Er$^{2+}$, Tm$^{3+}$), (Tm$^{2+}$, Yb$^{3+}$), (Yb$^{2+}$, Lu$^{3+}$).\\
If the value of set I is `x' and set II is `y', then find $|x-y|$.
\partsec{3}{MATHEMATICS}
\sect{SECTION-I(i)}
\qn{1} A drawer contains a mixture of red socks and blue socks, at most 2024 in all. It so happends that, when two socks are selected randomly without replacement, there is a probability of exactly $\dfrac12$ that both are red or both are blue. What is the largest possible number of red socks in the drawer that is consistent with this data ?
\begin{qoptions}
\opt{A}{1000}
\opt{B}{1012}
\opt{C}{990}
\opt{D}{946}
\end{qoptions}
\qn{2} Let $\dfrac{x-1}{2}=\dfrac{y-1}{3}=\dfrac{z+3}{2}$ and $\dfrac{x-k}{1}=\dfrac{y+1}{5}=\dfrac{z}{-1}$ be two lines intersect at $(\alpha,\beta,\gamma)$ then $\alpha+\beta+\gamma-k$ equals
\begin{qoptions}
\opt{A}{1}
\opt{B}{4}
\opt{C}{6}
\opt{D}{9}
\end{qoptions}
\qn{3} As shown in the diagram, the volume of tetrahedron $DABC$ is $\dfrac16$. Also $\angle ACB=45^\circ$ and $AD+BC+\dfrac{AC}{\sqrt2}=3$, then
\qfig{p1_math_s1_q03_1}
\begin{qoptions}
\opt{A}{$CD=\sqrt3$}
\opt{B}{$AD=\sqrt2$}
\opt{C}{$BC=2$}
\opt{D}{$AC=2$}
\end{qoptions}
\qn{4} Position vector of center of circle $|\vec r|=5$ and $\vec r\cdot(\hat i+\hat j+\hat k)=3\sqrt3$ is
\begin{qoptions}
\opt{A}{$\hat i+\hat j+\hat k$}
\opt{B}{$3(\hat i+\hat j+\hat k)$}
\opt{C}{$\sqrt3(\hat i+\hat j+\hat k)$}
\opt{D}{$\vec 0$}
\end{qoptions}
\sect{SECTION-I(ii)}
\qn{5} Let $\vec\alpha$ is the only point of intersection of three plane $\vec r\cdot\vec a=p_1$, $\vec r\cdot\vec b=p_2$ and $\vec r\cdot\vec c=p_3$ then which of following is/are correct
\begin{qoptions}
\opt{A}{$[\vec a\ \vec b\ \vec c]=0$}
\opt{B}{$[\vec a\ \vec b\ \vec c]\ne0$}
\opt{C}{$\vec\alpha=\dfrac{p_3(\vec a\times\vec b)+p_1(\vec b\times\vec c)+p_2(\vec c\times\vec a)}{[\vec a\ \vec b\ \vec c]}$}
\opt{D}{$\vec\alpha=p_3(\vec a\times\vec b)+p_1(\vec b\times\vec c)+p_2(\vec c\times\vec a)$}
\end{qoptions}
\qn{6} Three circles with radii 1,2 and 3 touch each other externally at point $A$, $B$ and $C$, then
\begin{qoptions}
\opt{A}{Area of $\Delta ABC=6/5$}
\opt{B}{Area of $\Delta ABC=7/5$}
\opt{C}{Circum-radius of $\Delta ABC=\dfrac54$}
\opt{D}{Circum-radius of $\Delta ABC=1$}
\end{qoptions}
\qn{7} Let $A$,$B$ and $C$ be three pair-wise independent events such that $P(A\cup B\cup C)=1$. Then which of the following is/are must be correct ?
\begin{qoptions}
\opt{A}{$1\le P(A)+P(B)+P(C)\le3$}
\opt{B}{$P(\bar A\cap\bar B\cap\bar C)=0$}
\opt{C}{$P(A)=P(B)=P(C)=1$}
\opt{D}{$P(A\cap B\cap C)=P(A)\cdot P(B)\cdot P(C)$}
\end{qoptions}
\qn{8} If $\vec a=\hat i+\hat j+\hat k$, $\vec b=\hat i-\hat j+\hat k$, $\vec c=2\hat i+3\hat j+\hat k$, $\vec c\times\vec b=\vec d\times\vec b$ and $\vec d\cdot\vec a=8$ then which of the following is/are correct ?
\begin{qoptions}
\opt{A}{$\vec d\cdot\vec b=4$}
\opt{B}{$\vec d\cdot\vec c=14$}
\opt{C}{$[\vec b\ \vec c\ \vec d]=0$}
\opt{D}{$[\vec a\ \vec b\ \vec c]=2$}
\end{qoptions}
\qn{9} If $A$ and $B$ are independent events such that the probability that they both occur simultaneously is $\dfrac18$ and the probability that neither of them will occurs is $\dfrac38$, then ($A$ is more probable then $B$)
\begin{qoptions}
\opt{A}{the probability that event $A$ will occur is $\dfrac12$}
\opt{B}{the probability that event $B$ will occur is $\dfrac14$}
\opt{C}{the probability that event $A$ will occur is less than $\dfrac34$}
\opt{D}{the probability that event $B$ will occur is more than $\dfrac13$}
\end{qoptions}
\qn{10} Suppose you play a game by tossing a fair coin repeatedly and independently. If a head comes up, you win Rs.1 and if a tail comes up you lose Rs.1. Suppose you start with Rs. $n$. Let $P(m,n)$ be the probability that you will get to Rs. $m$ first without running out of money ($m,n\in N$, $m>n$). Which of the following is/are CORRECT ?
\begin{qoptions}
\opt{A}{$P(51,50)=\dfrac12$}
\opt{B}{$P(325,52)=\dfrac4{25}$}
\opt{C}{$P(100,50)\ge\dfrac12$}
\opt{D}{$P(325,52)=\dfrac7{25}$}
\end{qoptions}
\sect{SECTION-II(i)}
\paraq{1}{2}
In $\Delta ABC$, $AB=AC$. Point $D$ and $E$ lie on $BC$ such that $BD=DC$ and $BE>CE$. $\tan\angle EAC$, $\tan\angle EAD$ and $\tan\angle EAB$ are in G.P.
\qn{1} $\angle EAD$ is equal to $\dfrac\pi k$, then the value of $k$ is
\qn{2} If $AE=10$ then $AD$ equals $\lambda\sqrt2$, then the value of $\lambda$ is
\paraq{3}{4}
Let a line with direction ratio 2,2,1 interset the line $\dfrac{x-7}3=\dfrac{y-5}2=\dfrac{z-3}1$ and $\dfrac{x-1}2=\dfrac{y+1}4=\dfrac{z+1}3$ at $A$ and $B$ respectively
\qn{3} If $(x_1,y_1,z_1)$ is $A$'s coordinate. Then find $x_1+y_1+z_1=$
\qn{4} Distance $AB$ is equal to
\paraq{5}{6}
Alice is in a wonderland of tunnels show below. At each intersection ($a$,$b$,$c$,$d$) he chooses a direction randomly (including possibly turning around). Once Alice reaches an exit, she will not return into the tunnels.
\qfig{p1_math_s3_para5_6_1}
Let $P(x_y)$ = Probability of exiting at $Y$ when currently at the intersection $x$.\\
Find
\qn{5} $15P(A_a)$ is equal to
\qn{6} $15P(b_A)$ is equal to
\sect{SECTION-II(ii)}
\qn{7} Let $A$ and $B$ be points $(2,0,0)$ and $(0,4,0)$ respectively. If $E$ is mid-point of $AB$ and $F$ is mirror image of origin with respect edge $AB$ and let $D$ the point on the line $\dfrac{x-1}0=\dfrac{y-2}0=\dfrac z1$ at 6 unit distance from the edge $AB$, then volume of tetrahedron $DFAB$ is\underline{\hspace{1.2cm}}
\qn{8} Let $O$--$ABC$ be a regular tetrahedron ($A(\vec a)$, $B(\vec b)$, $C(\vec c)$. $0$ is origin) with edge length 1.\\
Find $4\left|\vec a\times\vec b+\vec b\times\vec c+\vec c\times\vec a\right|^2=$?
\qn{9} From a bag $A$ containing 3 white and 5 black balls, 4 balls are transferred into an empty bag $B$. From bag $B$ 2 balls are taken and they both happen to be white. These balls are returned to the bag $B$. Now a ball is taken from bag $B$ is white; the probability of this event can be expressed as $p/q$ where $p$ and $q$ are relatively prime positive integers. Find $(q-p)$.
\papersec{2}
\partsec{1}{PHYSICS}
\sect{SECTION-I(i)}
\qn{1} $I$ is the image of a point object $O$ formed by spherical mirror then which of the following statement(s) is/are CORRECT?
\begin{qoptions}
\opt{A}{If $O$ and $I$ are on same side of the principal axis, then they have to be on opposite side of the mirror}
\opt{B}{If $O$ and $I$ are on opposite sides of the principal axis, then they have to be on same side of the mirror}
\opt{C}{If $O$ and $I$ are on opposite side of the principal axis, then they can be on opposite side of the mirror as well.}
\opt{D}{If $O$ and $I$ are on same side of the principal axis, then they have to be on same side of the mirror}
\end{qoptions}
\qn{2} A solid transparent cylinder is made of a material whose refractive index depends on the distance $r$ from its axis of symmetry according to the formula
\[n=1+(n_0-1)\cdot\Bigl(1-\Bigl(\frac r{r_0}\Bigr)\Bigr)\]
Here $r_0$ is the radius of the cylinder, $n_0=5/4$. A narrow laser beam falls perpendicular to the end of the cylinder in the direction parallel to the axis of the cylinder on a point located at a distance of $r_0$ from the axis of cylinder. Choose the correct statement(s) relating to the path followed by the narrow beam within the cylinder
\qfig{p2_phy_s1_q02_1}
\begin{qoptions}
\opt{A}{The beam follows a helical path around the axis of the cylinder}
\opt{B}{The beam follows a path in a plane containing the axis of the cylinder and its initial direction of propagation}
\opt{C}{The path crosses the axis at a distance of $5r_0\ln3$ from the end of which beam is incident}
\opt{D}{The path crosses the axis at a distance of $4r_0\,\ell n2$ from the end on which the beam is incident}
\end{qoptions}
\qfullbegin
\qn{3} A glass prism is immersed in a hypothetical liquid as shown in figure (i). The curves showing the refractive index $n$ as a function of wavelength $\lambda$ for glass and liquid. When a ray of white light is incident on the prism parallel to the base
\qfig{p2_phy_s1_q03_1}
\qfullsplit
\begin{qoptions}
\opt{A}{yellow ray travels without deviation}
\opt{B}{blue ray is deviated towards the vertex}
\opt{C}{red ray is deviated towards the base}
\opt{D}{there is no dispersion}
\end{qoptions}
\qfullend
\qn{4} Figure shows converging rays on a glass slab placed in air. $O$ is point of intersection of incident rays. $I_1$ and $I_2$ are image formed after refraction at $S_1$ and $S_2$ surface respectively. Mark the correct option(s)
\qfig{p2_phy_s1_q04_1}
\begin{qoptions}
\opt{A}{for refraction at $S_2$, $I_1$ is real object and $I_2$ is real image}
\opt{B}{for refraction at $S_1$, $O$ is real object and $I_1$ is its real image.}
\opt{C}{for refraction at surface $S_1$, $O$ is virtual object and $I_1$ is real image}
\opt{D}{for refraction at surface $S_2$, $I_1$ is virtual object and $I_2$ is real image}
\end{qoptions}
\qn{5} The two faces of a prism of side length $a$ are silvered. A ray incidents normally unsilvered face as shown. Then
\qfig{p2_phy_s1_q05_1}
\begin{qoptions}
\opt{A}{net deviation of ray through prism is $90^\circ$}
\opt{B}{net deviation of ray through prism is $180^\circ$}
\opt{C}{length of path travelled by light ray inside the prism is $\sqrt5\,a$}
\opt{D}{length of path travelled by light ray inside the prism is $\sqrt3\,a$}
\end{qoptions}
\qn{6} A light ray is incident on a boundary as shown.
\qfig{p2_phy_s1_q06_1}
\begin{qoptions}
\opt{A}{If $\mu_2<\sqrt3$, light will enter medium-3 with deviation $30^\circ$.}
\opt{B}{If $\mu_2>\sqrt3$, light will enter medium-3 with deviation $30^\circ$.}
\opt{C}{If $\mu_2>\sqrt3$, light will never return to medium-1 again.}
\opt{D}{If $\mu_2<\sqrt3$, light will again return to medium-1.}
\end{qoptions}
\sect{SECTION-I(ii)}
\paraq{7}{8}
Whenever an object is placed between the two mirrors at an angle $\theta$ on its bisector. It forms several images which lies on a circle of radius equal to the distance between the vertex of mirror and the object.
\qfig{p2_phy_s2_para7_8_1}
\qn{7} Draw graph between the area formed by joining all the images and object and the distance between the vertex and object
\begin{qoptions}
\opt{A}{\qfigi{p2_phy_s2_q07_A}}
\opt{B}{\qfigi{p2_phy_s2_q07_B}}
\opt{C}{\qfigi{p2_phy_s2_q07_C}}
\opt{D}{\qfigi{p2_phy_s2_q07_D}}
\end{qoptions}
\qn{8} If the angle between the two mirror is $60^\circ$. Find the area of regular polygon bounded by the images and the object
\begin{qoptions}
\opt{A}{$\dfrac{x^2\left(4\sqrt3+1\right)}{4}$}
\opt{B}{$x^2\sqrt{\left(\dfrac{4\sqrt3+1}{4}\right)x}$}
\opt{C}{$\dfrac{3\sqrt3x^2}{2}$}
\opt{D}{$\dfrac{5\sqrt3x^2}{4}$}
\end{qoptions}
\paraq{9}{10}
There is a glass rod placed in the air whose refractive index $n=1.5$, the central axis has length $L=45$ cm, one end is a convex spherical surface with a radius of $R_1=10$ cm. This surface faces object. We want to make the function of the glass rod equivalent to that of an ideal astronomical telescope (magnifying the object at infinity on the main optical axis so that the final image is formed at infinity on the main optical axis). Take the central axis of the glass rod as the main optical axis.
\qfig{p2_phy_s2_para9_10_1}
\qn{9} What kind of spherical surface should the other end be?
\begin{qoptions}
\opt{A}{Concave with radius 5 cm.}
\opt{B}{Convex with radius 5 cm.}
\opt{C}{Concave with radius 10 cm.}
\opt{D}{Convex with radius 20 cm.}
\end{qoptions}
\qn{10} What is the angular magnification of this system for the object at infinity?
\begin{qoptions}
\opt{A}{4}
\opt{B}{2}
\opt{C}{3}
\opt{D}{6}
\end{qoptions}
\sect{SECTION-II(i)}
\paraq{1}{2}
There are two spheres of transparent materials. Sphere $A$ has refractive index of $\mu=\dfrac85$ and refractive index of sphere $B$ is not known to us (consider it as $\mu_2$). Radii of each sphere is 50 cm. Both the spheres are cut so that the section of cut is 30 cm from center. Later the smaller parts struck together as shown.
\qfig{p2_phy_s3_para1_2_1}
A narrow laser beam from point $S$ (source) is incident on the curved surface of the part of sphere $A$ at point $P$. The length of beam from $S$ to $P$ is 50 cm. Also the angle made by beam with the principal axis is $\theta=37^\circ\ \bigl(\sin\theta=\dfrac35\bigr)$. The beam finally emerges from the curved surface of the part of sphere $B$. The beam follows a total deviation of $53^\circ$.
\qn{1} If the image of first refraction is formed at a distance $D$ cm from the surface of $A$, find the value of $\dfrac1D$.
\qn{2} If the beam after emerging from $B$ meets the principal axis at $I$, the total length of beam (in cm) from $S$ to $I$ will be given by $\dfrac{10}7Q$. Find the value of $Q$.
\paraq{3}{4}
For a glass prism the plot of deviation ($\delta$) vs angle of incidence ($i$) is as shown. If the refractive index of the glass is $\mu$ and value of angle is $i_1$. Then :
\qfig{p2_phy_s3_para3_4_1}
\qn{3} Find the value of $i_1$ (in degree) :
\qn{4} Find the refractive index of glass :
\paraq{5}{6}
As shown in the figure, there is large vessel filled with a liquid with a refractive index $n$, and a symmetrical cone-shaped thin-walled transparent cover $ADB$ is stuck in the middle of the vessel. The top angle is $2\phi$, inside of the cone is filled with air, the whole cover is submerged in the liquid, and there is a bright spot $C$ at the midpoint of $AB$ at the bottom of the tank. Take $n=5/3$.
\qfig{p2_phy_s3_para5_6_1}
\qn{5} What is the minimum value of $\tan\phi$ so that an observer outside may be able to see the light coming from $C$?
\qn{6} If $\phi=87^\circ$, what will be the angle (in degree) between the extreme rays perceived by an observer just above $D$?
\sect{SECTION-II(ii)}
\qn{7} A ray of light is incident on a prism as shown in the figure. The total deviation suffered by the ray is $\dfrac{x^\circ}2$. Then find $x$. ($\angle BAC=1^\circ$ and $\angle CAD=2^\circ$, refractive indices are shown in figure.)
\qfig{p2_phy_s4_q07_1}
\qn{8} On the surface of a liquid reservoir of depth 3m is a circular raft of radius 6m. The surface is illuminated by light incident from all the directions. Determine the radius (in m) of the shadow of the raft at the bottom of the reservoir. The refractive index of the liquid is 1.25.
\qn{9} A point object is placed at a distance of 25 cm from a convex lens of focal length 20 cm. If a glass slab of thickness $t$ and refractive index 1.5 is inserted between the lens and object. The image is formed at infinity. Find the thickness $t$ (in cm)?
\partsec{2}{CHEMISTRY}
\sect{SECTION-I(i)}
\qn{1} Which of the following option(s) is/are correct for the structure of H$_2$O$_2$ in gas phase and solid phase?
\qfig{p2_chem_s1_q01_1}
\begin{qoptions}
\opt{A}{$X>Y$ (O -- O) bond length}
\opt{B}{$X<Y$ (O -- O) bond length}
\opt{C}{$\alpha>\beta$ (dihedral angle)}
\opt{D}{$\alpha<\beta$ (dihedral angle)}
\end{qoptions}
\qn{2} Which of the following property is not same of both H$_2$ and D$_2$?
\begin{qoptions}
\opt{A}{Enthalpy of fusion}
\opt{B}{Enthalpy of vaporization}
\opt{C}{Internuclear distance}
\opt{D}{Enthalpy of bond dissociation}
\end{qoptions}
\qn{3} The used organic solvent in industrial preparation of H$_2$O$_2$ must
\begin{qoptions}
\opt{A}{Dissolve quinol}
\opt{B}{Dissolve quinone}
\opt{C}{Resist oxidation}
\opt{D}{Be immiscible with water}
\end{qoptions}
\qn{4} Lanthanoids (Ln) when,
\begin{qoptions}
\opt{A}{Heated with S, produces Ln$_2$S$_3$}
\opt{B}{Burns in O$_2$, produces Ln$_2$O$_3$}
\opt{C}{React with H$_2$O produces Ln(OH)$_3$ + H$_2$}
\opt{D}{React with halogens, produces LnX$_3$}
\end{qoptions}
\qn{5} Select the correct matching of maximum prescribed concentration in drinking water with corresponding metals.
\begin{qoptions}
\opt{A}{Fe -- 0.2 mg dm$^{-3}$}
\opt{B}{Cu -- 0.05 mg dm$^{-3}$}
\opt{C}{Zn -- 5 mg dm$^{-3}$}
\opt{D}{Cd -- 0.2 mg dm$^{-3}$}
\end{qoptions}
\qn{6} The pair(s) of complexes where in both exhibit tetrahedral geometry is(are)
\begin{qoptions}
\opt{A}{[FeCl$_4$]$^-$ and [Fe(CO)$_4$]$^{2-}$}
\opt{B}{[Co(CO)$_4$]$^-$ and [CoCl$_4$]$^{2-}$}
\opt{C}{[Ni(CO)$_4$] and [Ni(CN)$_4$]$^{2-}$}
\opt{D}{[Cu(py)$_4$]$^+$ and [Cu(CN)$_4$]$^{3-}$}
\end{qoptions}
\sect{SECTION-I(ii)}
\paraq{7}{8}
Structural isomerism arises due to the difference in the type of chemical linkages and distribution of ligands within and outside the coordination sphere whereas stereoisomers differ only with respect to the spatial orientation of ligands in space around the metal ion.
\qn{7} Select the type of complex which have the ratio of total number of enentiomeric pairs to the total number of stereoisomers equal to $\dfrac34$.
\begin{qoptions}
\opt{A}{[M(AB)a$_2$bc]}
\opt{B}{[M(AB)$_2$a$_2$]}
\opt{C}{[Ma$_2$b$_2$cd]}
\opt{D}{[M(AB)$_2$ab]}
\end{qoptions}
\qn{8} In the complex [Pt(Cl)(NO$_2$)(NO$_3$)(SCN)]$^{2-}$,\\
Total number of possible isomers $=a$\\
Total number of unpaired electrons on central metal ion $=b$\\
Number of inner/outer d-orbitals involved in the hybridisation of the complex $=c$\\
The value of $(a\times c)+b$ is :-
\begin{qoptions}
\opt{A}{12}
\opt{B}{13}
\opt{C}{3}
\opt{D}{0}
\end{qoptions}
\paraq{9}{10}
Co-ordination compounds are used as catalyst for many industrial processes. Example include Rhodium complex used for the hydrogenation of alkenes, this catalyst is called as Wilkinson catalyst.
\qn{9} What is true about Wilkinson catalyst?
\begin{qoptions}
\opt{A}{Co-ordination number of Rhodium is three}
\opt{B}{Oxidation number of Rhodium is $+2$}
\opt{C}{Co-ordination number of Rhodium is six}
\opt{D}{Oxidation number of Rhodium is $+1$}
\end{qoptions}
\qn{10} There are four different compounds of Co$^{3+}$ with ligand Cl$^-$, NH$_3$, H$_2$O and CN$^-$ that are A, B, C and D. Consider the following information.
\begin{center}\small
\begin{tabular}{|c|c|}
\hline Compound & Wavelength of light absorbed (nm) \\\hline
A & 535 \\\hline
B & 500 \\\hline
C & 475 \\\hline
D & 310 \\\hline
\end{tabular}
\end{center}
The possible combination of A, B, C and D
\begin{qoptions}
\opt{A}{A$\to$\allowbreak[CoCl(NH$_3$)$_5$]$^{2+}$, B$\to$\allowbreak[Co(NH$_3$)$_5$\allowbreak(H$_2$O)]$^{3+}$, C$\to$\allowbreak[Co(NH$_3$)$_6$]$^{3+}$, D$\to$\allowbreak[Co(CN)$_6$]$^{3-}$}
\opt{B}{A$\to$\allowbreak[Co(NH$_3$)$_5$\allowbreak(H$_2$O)]$^{3+}$, B$\to$\allowbreak[CoCl(NH$_3$)$_5$]$^{2+}$, C$\to$\allowbreak[Co(NH$_3$)$_6$]$^{3+}$, D$\to$\allowbreak[Co(CN)$_6$]$^{3-}$}
\opt{C}{A$\to$\allowbreak[Co(CN)$_6$]$^{3-}$, B$\to$\allowbreak[Co(NH$_3$)$_6$]$^{3+}$, C$\to$\allowbreak[Co(NH$_3$)$_5$\allowbreak(H$_2$O)]$^{3+}$, D$\to$\allowbreak[CoCl(NH$_3$)$_5$]$^{2+}$}
\opt{D}{A$\to$\allowbreak[CoCl(NH$_3$)$_5$]$^{2+}$, B$\to$\allowbreak[Co(CN)$_6$]$^{3-}$, C$\to$\allowbreak[Co(NH$_3$)$_6$]$^{3+}$, D$\to$\allowbreak[Co(NH$_3$)$_5$\allowbreak(H$_2$O)]$^{3+}$}
\end{qoptions}
\sect{SECTION-II(i)}
\paraq{1}{2}
All the lanthanoids are silvery white soft metals and tarnish rapidly in air. They have typical metallic structure and are good conductors of heat and electricity.
\qn{1} The atomic number of element having unusually large melting point than the other lanthanides is :-
\qn{2} Consider the electronic configuration of lanthanides series. If the unpaired electrons in the ground state of Ce$^{+3}$, Yb$^{+3}$ and Gd$^{+3}$ are $a$, $b$ and $c$ respectively, then the value of $a+b+c$ ?
\paraq{3}{4}
Hydrogen peroxide is an important chemical used in pollution control treatment of domestic and industrial effluents.
\qn{3} Consider the following species :-\\
Fe$^{2+}$, PbS, MnO$_4^-$, HOCl, Mn$^{2+}$, I$_2$\\
Number of species which shows reducing action with H$_2$O$_2$ in acidic medium with their valence factor `2' $=a$\\
Number of species which shows oxidizing action with H$_2$O$_2$ in acidic medium with their valence factor `8' $=b$\\
Number of species which shows reducing action with H$_2$O$_2$ in basic medium with their valence factor `3' $=c$\\
Calculate $\dfrac{b}{a+c}$.
\qn{4} If commercially marketed sample of H$_2$O$_2$ solution is 4.54 V, then \% H$_2$O$_2$ in the solution is :-
\paraq{5}{6}
In 1980's atmospheric scientists working in Antarctica reported about depletion of ozone layer commonly known as ozone hole over the south pole. It was found that a unique set of conditions was responsible for the ozone hole.
\qn{5} How many of the following species (free radicals and molecules) are responsible for the depletion of ozone layer?\\
CF$_2$Cl$_2$, ClONO$_2$, CH$_4$, ClO, CH$_3$, Cl, CH$_2$O, HOCl, NO$_2$, HCl
\qn{6} When chlorine free radical reacts with methane in summer season, the free radical product formed have total `x' number of hybrid orbitals. Find `x'.
\sect{SECTION-II(ii)}
\qn{7} A hydrated isomer of CrCl$_3$.6H$_2$O, for which no geometrical isomer is possible, when 0.1 mole of such complex in 1 L solution is treated with excess of AgNO$_3$, mass of white ppt. formed is :-\\
(Give nearest rounded off integer)
\qn{8} In Mischmetall, ratio of (approx.) \% Lanthanoid metal used to (approx.) \% iron used is :-
\qn{9} Find out the number of following orders which are INCORRECT against the mentioned properties :\\
(i) Ortho hydrogen $>$ Para hydrogen (Stability at low temperature)\\
(ii) Ethanol $>$ Glycerol (Viscosity)\\
(iii) D$_2$ $<$ He (Boiling Point)\\
(iv) HF $>$ H$_2$O (Melting point)\\
(v) H$_3$BO$_3$ $>$ BF$_3$ (Melting point)\\
(vi) NH$_3$ $<$ SbH$_3$ (Boiling point)\\
(vii) H$_2$O$_2$ $>$ H$_2$O (Strength of hydrogen bond)\\
(viii) H$_2$O $<$ D$_2$O (Freezing point)\\
(ix) HF$_2^-$ $<$ HF (Strength of hydrogen bond)\\
(x) D$_2$O $<$ H$_2$O (Bond energy)
\partsec{3}{MATHEMATICS}
\sect{SECTION-I(i)}
\qn{1} Letters of the word ``UNIVERSITY'' are arranged randomly. Then
\begin{qoptions}
\opt{A}{Probability that both `I' do not come together is $4/5$}
\opt{B}{Probability that all vowels are seperated is $1/6$}
\opt{C}{Probability that all vowels are together is $1/30$}
\opt{D}{Probability that word starts and ends with a consonant is $1/3$}
\end{qoptions}
\qn{2} Let $\vec p=\vec a\times(\vec b+\vec c)$, $\vec q=\vec b\times(\vec c+\vec a)$ and $\vec r=\vec c\times(\vec a+\vec b)$. Then
\begin{qoptions}
\opt{A}{$\vec p,\vec q,\vec r$ are coplanar}
\opt{B}{If $x\vec p+y\vec q+z\vec r=\vec 0$; $x,y,z\in R$ then $\dfrac{x+z}y=2$}
\opt{C}{$\vec p\cdot\vec q+\vec q\cdot\vec r+\vec r\cdot\vec p=-\bigl(|\vec p|^2+|\vec q|^2+|\vec r|^2\bigr)$}
\opt{D}{$\vec p\cdot\vec r=\dfrac12\bigl(|\vec q|^2-|\vec p|^2-|\vec r|^2\bigr)$}
\end{qoptions}
\qn{3} Sides $a$, $b$, $c$ of $\Delta ABC$ are roots of the equation $3x^3-px^2+qx-4=0$ ($p,q\in R$) such that $\dfrac a3+\dfrac b4+\dfrac c3=1$. Then
\begin{qoptions}
\opt{A}{$\Delta ABC$ is equilateral}
\opt{B}{$\Delta ABC$ is isosceles}
\opt{C}{Area of $\Delta ABC=\dfrac{2\sqrt5}9$}
\opt{D}{Circum radius of $\Delta ABC=\dfrac3{10}$}
\end{qoptions}
\qn{4} A ray is sent along the line $\dfrac x2=\dfrac{y-2}2=\dfrac{z-1}0$ and is reflected by the plane $x=0$ at point $P$. The reflected ray is again reflected by the plane $x+2y=0$ at point $Q$. The initial ray and final reflected ray meet at point $R$. Then
\begin{qoptions}
\opt{A}{The co-ordinates of $Q$ is $(4,-2,1)$}
\opt{B}{The co-ordinates of $R$ is $(-3,-1,1)$}
\opt{C}{The centroid of $\Delta PQR$ is $(0,0,0)$}
\opt{D}{The centroid of $\Delta PQR$ is $\Bigl(\dfrac13,\dfrac{-1}3,1\Bigr)$}
\end{qoptions}
\qn{5} A bag contains some white and some black balls, all combinations being equally likely. The total number of balls in the bag is 12. Four balls are drawn from the bag without replacement.
\begin{qoptions}
\opt{A}{Probability that all balls are black is $1/5$}
\opt{B}{Probability that all balls are black is $1/3$}
\opt{C}{Probability that two balls are black and two are white is $1/4$}
\opt{D}{Probability that two balls are black and two are white is $1/5$}
\end{qoptions}
\qn{6} Which of the following statement is/are true ?
\begin{qoptions}
\opt{A}{If $\vec a,\vec b,\vec c$ are co-terminus edges of a regular tetrahedron then $\vec a\cdot\vec b=\vec b\cdot\vec c=\vec c\cdot\vec a$}
\opt{B}{If $\vec a,\vec b$ are two unit vector inclined at $\dfrac\pi3$, then $[\vec a\ \ \vec b+\vec a\times\vec b\ \ \vec b]=\dfrac34$}
\opt{C}{Three points having position vectors $\vec a,\vec b$ and $\vec c$ are collinear if $[\vec a\ \ \vec b\ \ \vec c]=0$}
\opt{D}{If $\vec a\times\vec b=\vec c$ and $\vec b\times\vec c=\vec a$ then $\vec a\cdot\vec b=0$}
\end{qoptions}
\sect{SECTION-I(ii)}
\paraq{7}{8}
Let $B_n$ denotes the event that $n$ fair dice are rolled once with $P(B_n)=\dfrac1{2^n}$; $n\in N$. Hence $B_1,B_2,B_3,\ldots,B_n$ are pairwise mutually exclusive events as $n\to\infty$. The event $E$ occurs with atleast one of the event $B_1,B_2,\ldots,B_n$ and denotes that the sum of numbers appearing on dice is $S$.
\qn{7} Probability that greatest number on dice is 4, if three dice are known to have been rolled, is :
\begin{qoptions}
\opt{A}{$\dfrac{37}{216}$}
\opt{B}{$\dfrac{64}{216}$}
\opt{C}{$\dfrac{27}{216}$}
\opt{D}{$\dfrac{31}{216}$}
\end{qoptions}
\qn{8} If $S=3$, then $P\Bigl(\dfrac{B_2}S\Bigr)$ is equal to -
\begin{qoptions}
\opt{A}{$\dfrac8{169}$}
\opt{B}{$\dfrac{24}{169}$}
\opt{C}{$\dfrac{72}{169}$}
\opt{D}{$\dfrac{16}{169}$}
\end{qoptions}
\paraq{9}{10}
Lines $\dfrac{x-3}2=\dfrac{y-2}3=\dfrac{z-1}\lambda$ and $\dfrac{x-2}3=\dfrac{y-3}2=\dfrac{z-2}3$ intersect at point $P$. Then -
\qn{9} Equation of the plane containing both the lines is
\begin{qoptions}
\opt{A}{$x+5y-3z=10$}
\opt{B}{$x+6y+5z=20$}
\opt{C}{$x+6y-5z=10$}
\opt{D}{$x+5y+3z=20$}
\end{qoptions}
\qn{10} The distance of point $P$ from the line $\dfrac{x-1}2=\dfrac{y-3}4=\dfrac{z-2}3$ measured parallel to the plane $3x+2y-2z+15=0$ is -
\begin{qoptions}
\opt{A}{$\dfrac{17}4$}
\opt{B}{$\dfrac{3\sqrt{21}}4$}
\opt{C}{$\dfrac{3\sqrt{21}}2$}
\opt{D}{$\dfrac{17}2$}
\end{qoptions}
\sect{SECTION-II(i)}
\paraq{1}{2}
On a normal standard die one of the 21 dots from any one of the six faces is removed at random with each dot equally likely to be chosen. The die is then rolled.
\qn{1} Probability of die showing even number of dots is $\dfrac ab$, where $a$ \& $b$ are co-prime, then value of $a+b$ is
\qn{2} Probability of die showing odd number of dots is $\dfrac pq$, where $p$ and $q$ are co-prime, then value of $q^2-p^2$ is -
\paraq{3}{4}
In a $\Delta ABC$, the incircle touches sides $BC$, $CA$ and $AB$ at $D$, $E$, $F$ respectively such that $AF$, $BD$, $CE$ are consecutive even integers and inradius is $\dfrac{4\sqrt6}3$.
\qn{3} If $R$ is circum radius then find value of $\dfrac{24}{25}R^2$.
\qn{4} If $\ell_1,\ell_2,\ell_3$ are length of medians then find value of $\ell_1^2+\ell_2^2+\ell_3^2$.
\paraq{5}{6}
Let $\vec a,\vec b$ and $\vec c$ be three non-coplanar unit vector such that the angle between every pair of them is $\dfrac\pi3$. If $\vec a\times\vec b+\vec b\times\vec c=p\vec a+q\vec b+r\vec c$, where $p,q,r$ are scalars, then
\qn{5} The value of $[\vec a\ \ \vec b\ \ \vec c]^4$ is \underline{\hspace{1.5cm}}
\qn{6} The value of $20p^2+21q^2-22r^2$ is \underline{\hspace{1.5cm}}
\sect{SECTION-II(ii)}
\qn{7} Let position vectors of vertices of $\Delta ABC$ be $3\hat i+4\hat j+5\hat k$, $-2\hat i+6\hat j+\sqrt{10}\hat k$, $5\hat j-5\hat k$. Let $DEF$ be triangle formed by joining the mid points of sides of $\Delta ABC$. If $\vec P$ is the position vector of orthocentre of $\Delta DEF$ and position vector of $Q$ is $2\hat i-3\hat j+6\hat k$, then find $\left|\overrightarrow{PQ}\right|$.
\qn{8} $S$ is a set containing 5 elements. Two subsets $A$ and $B$ of set $S$ are selected at random. The probability that $A\cap B$ contains 2 elements is $P$, then find value of $(512)P$
\qn{9} Let $\vec a,\vec b,\vec c$ are three non-zero, non-collinear vectors and $\vec r_1=p\vec a+q\vec b+\vec c$ and $\vec r_2=\vec a+p\vec b+q\vec c$. If the vectors $\vec r_1+2\vec r_2$ and $2\vec r_1+\vec r_2$ are collinear then find value of $p^2+q^2$.

\end{multicols}
\end{document}
